\chapter{Melasma (Chloasma)}

\section*{Definition and Overview}
Melasma, historically referred to as chloasma or ``the mask of pregnancy,'' is a common, chronic, acquired cutaneous disorder characterised by symmetrical, hyperpigmented macules and patches on sun-exposed areas of the body. The pigmentation most frequently affects the face, particularly the cheeks, forehead, upper lip, nose, and chin. The condition is caused by hyperfunctional melanocytes producing excessive amounts of the pigment melanin, which is deposited within the epidermis, dermis, or both. Although melasma is entirely benign and carries no risk of malignant transformation, it represents a significant aesthetic concern that can cause profound psychosocial distress, anxiety, and impaired quality of life. The condition is strongly linked to ultraviolet (UV) radiation, visible light, genetic predisposition, and hormonal fluctuations.

\section*{Epidemiology}
Melasma is a highly prevalent condition worldwide, particularly in individuals with darker skin temperaments (Fitzpatrick skin phototypes III to VI), such as those of Asian, Hispanic, Middle Eastern, and African descent. It is far more common in females than in males, with women accounting for approximately 90\% of all cases. The onset typically occurs during the reproductive years, between the ages of 20 and 50. In Australia, the combination of high ambient solar radiation, outdoor lifestyle patterns, and a diverse population contributes to a high clinical incidence. The pigmentation is characteristically dynamic, typically darkening during the summer months due to increased UV exposure and fading or improving during the winter.

\section*{Aetiology and Pathophysiology}
The pathogenesis of melasma is complex and multifactorial, involving melanocyte hyperactivity stimulated by external and internal factors:
\begin{itemize}
    \item \textbf{Ultraviolet and visible light exposure:} UV radiation directly stimulates melanogenesis by upregulating melanocyte-stimulating hormone (MSH) and inducing reactive oxygen species. Visible light, particularly high-energy visible (HEV) blue light, also triggers pigment production in darker skin types.
    \item \textbf{Hormonal stimulation:} Elevated levels of oestrogen and progesterone sensitise melanocytes. This is evidenced by the frequent onset or worsening of melasma during pregnancy, the use of combined oral contraceptive pills, and hormone replacement therapy.
    \item \textbf{Genetic predisposition:} A positive family history is reported in over 50\% of affected individuals, suggesting a strong genetic susceptibility and familial clustering.
    \item \textbf{Thyroid autoimmunity:} An increased association has been observed between melasma and autoimmune thyroid disorders, such as Hashimoto's thyroiditis.
    \item \textbf{Chronic inflammation and vascularity:} Upregulation of vascular endothelial growth factor (VEGF) in melasma lesions leads to an increased number of prominent dermal blood vessels (telangiectasias), which interact with and stimulate overlying melanocytes.
    \item \textbf{Barrier dysfunction:} Subclinical cutaneous barrier impairment and accelerated cellular ageing (solar elastosis) within the dermal extracellular matrix due to chronic photodamage.
\end{itemize}

\section*{Clinical Features}
Melasma presents with distinct, symmetrical hyperpigmented lesions on the face and occasionally other sun-exposed sites:
\begin{itemize}
    \item \textbf{Symmetrical hyperpigmented patches:} Irregularly shaped, macular patches of light-to-dark brown, grey-brown, or slate-grey pigmentation.
    \item \textbf{Centrofacial distribution:} The most common pattern (occurring in 65\% of cases), involving the forehead, cheeks, nose, upper lip, and chin.
    \item \textbf{Malar distribution:} Pigmentation restricted to the cheeks and nose.
    \item \textbf{Mandibular distribution:} Pigmentation occurring along the ramus of the jawline.
    \item \textbf{Asymptomatic nature:} The patches are completely asymptomatic, with no associated pruritus, scaling, pain, or warmth.
    \item \textbf{Seasonal variation:} Distinct darkening of the lesions following brief sun exposure in summer, and slow fading during periods of low UV index.
\end{itemize}

\section*{Differential Diagnosis}
Because facial hyperpigmentation can arise from various inflammatory, metabolic, or neoplastic conditions, melasma must be distinguished from:
\begin{itemize}
    \item \textbf{Post-inflammatory hyperpigmentation (PIH):** Hyperpigmentation following a cutaneous injury, burn, or inflammatory dermatosis (e.g. acne, eczema), distinguished by a history of preceding lesions and a non-symmetrical distribution.
    \item \textbf{Solar lentigines:} Well-circumscribed, discrete brown macules (``age spots'') resulting from localized UV damage, which lack the large, confluent, symmetrical pattern of melasma.
    \item \textbf{Lichen planus pigmentosus (LPP):} An inflammatory condition presenting with slate-grey to dark brown macules in intertriginous or sun-exposed areas, often accompanied by mild pruritus and distinct histopathology.
    \item \textbf{Riehl's melanosis:} A contact-dermatitis-induced pigmentation, typically triggered by cosmetic ingredients, presenting with diffuse, reticular grey-brown pigmentation.
    \item \textbf{Drug-induced hyperpigmentation:} Symmetrical or asymmetrical slate-grey pigmentation caused by the systemic deposition of medications (such as minocycline, amiodarone, or antimalarials).
\end{itemize}

\section*{When to Seek Medical Care}
Patients should consult a general practitioner or dermatologist if they notice new or changing skin pigmentation, to obtain an accurate diagnosis and discuss cosmetic treatment options.

\textbf{Call 000 (or local emergency services) for:}
\begin{itemize}
    \item Signs of a severe allergic reaction (anaphylaxis) to prescribed topical treatments: difficulty breathing, swelling of the face, lips, tongue, or throat, and severe dizziness.
    \item Severe, widespread skin blistering, peeling, or intense burning pain following chemical peels or laser treatments.
    \item A new, rapidly changing, bleeding, or asymmetrical dark spot or mole (to rule out cutaneous melanoma).
\end{itemize}

\section*{Diagnosis and Investigations}
The diagnosis of melasma is primarily clinical, but specific tools help determine the depth of pigment and guide treatment:
\begin{itemize}
    \item \textbf{Clinical evaluation:} Inspection of the symmetry, border characteristics, and distribution of the facial patches.
    \item \textbf{Wood's lamp examination:} Inspection under long-wave ultraviolet light (365 nm) to determine the depth of melanin. Epidermal melasma shows enhanced contrast under the lamp; dermal melasma becomes less apparent; mixed melasma shows a combination of both.
    \item \textbf{Dermoscopy:} High-magnification skin inspection showing a characteristic brown pseudonetwork, follicular sparing, telangiectasias, and ruling out ochronosis (which shows dark blue-black banana-shaped structures).
    \item \textbf{Skin biopsy:} Rarely performed, but indicated if the diagnosis is uncertain to rule out conditions like melanoma or lichen planus pigmentosus.
    \item \textbf{Thyroid screening:} Evaluation of TSH and free thyroid hormones if clinical features suggest concurrent autoimmune thyroid disease.
\end{itemize}

\section*{Management}
Management requires a long-term, multimodal approach combining strict photoprotection, topical therapies, and lifestyle adjustments:
\begin{itemize}
    \item \textbf{Strict photoprotection:} The absolute foundation of therapy. Daily use of broad-spectrum SPF 50+ sunscreen containing physical blockers (zinc oxide and titanium dioxide) to reflect UV, and iron oxides to block visible blue light. This must be combined with wide-brimmed hats and sun-avoidance.
    \item \textbf{Topical tyrosinase inhibitors:} The gold standard is hydroquinone (2--4\%), which inhibits the conversion of dopa to melanin. It is often combined with tretinoin and a mild corticosteroid (triple combination cream or Kligman's formula) to enhance efficacy and reduce irritation.
    \item \textbf{Adjuvant topical agents:} Non-hydroquinone options including azelaic acid, kojic acid, ascorbic acid (vitamin C), and cysteamine cream.
    \item \textbf{Systemic therapy:} Low-dose oral tranexamic acid (typically 250 mg twice daily) for severe, refractory cases, which works by inhibiting plasminogen activation and reducing melanocyte-keratinocyte interactions.
    \item \textbf{Superficial chemical peels:} Glycolic acid or salicylic acid peels to accelerate epidermal cell turnover and remove pigment, performed with caution to avoid triggering post-inflammatory hyperpigmentation.
    \item \textbf{Laser and light therapies:} Low-fluence Q-switched Nd:YAG or picosecond lasers may be used for resistant dermal pigment, though there is a high risk of rebound hyperpigmentation.
\end{itemize}

\section*{Complications}
While melasma does not pose physical health risks, it is associated with significant local and psychological complications:
\begin{itemize}
    \item \textbf{Psychosocial distress:} Lowered self-esteem, self-consciousness, social anxiety, and depression due to the facial appearance.
    \item \textbf{Exogenous ochronosis:} A rare, disfiguring complication caused by long-term, uninterrupted use of high-concentration hydroquinone, presenting as soot-like bluish-black pigmentation.
    \item \textbf{Post-inflammatory hyperpigmentation (PIH):} Worsening of the pigmentation caused by irritation from aggressive lasers, strong chemical peels, or topical retinoids.
    \item \textbf{Contact dermatitis and skin atrophy:} Erythema, peeling, burning, and thinning of the skin from prolonged application of prescription topical treatments.
\end{itemize}

\section*{Prognosis and Outcomes}
Melasma is a chronic condition characterized by frequent relapses and resistance to treatment. While significant clearance or fading of the patches can be achieved using triple combination creams and sun protection, complete and permanent cure is rare. Epidermal melasma carries a much better prognosis and responds relatively quickly (within 2--3 months) to topical therapies. Dermal melasma is highly persistent and often resolves only partially. Melasma triggered by pregnancy or oral contraceptives may resolve slowly after delivery or cessation of the medication, but in many cases, it persists for years.

\section*{Prevention and Self-Care}
Ongoing self-care and preventative habits are crucial to control melasma and prevent relapse:
\begin{itemize}
    \item \textbf{Year-round sunscreen application:} Apply a broad-spectrum physical sunscreen every single day, regardless of weather, season, or indoor activity, and reapply every 2 hours when outdoors.
    \item \textbf{Use tinted sunscreens:} Select sunscreens containing iron oxides, which provide essential protection against the visible blue light emitted by both the sun and digital screens.
    \item \textbf{Physical barriers:} Wear broad-brimmed hats and sunglasses, and use car window UV-blocking films.
    \item \textbf{Avoid skin irritation:} Steer clear of harsh facial scrubs, alcohol-based cosmetics, or perfumed products that can cause subclinical inflammation.
    \item \textbf{Review medications:} Discuss alternative, non-hormonal contraceptive methods with your doctor if oral contraceptives are suspected of triggering your melasma.
\end{itemize}

\section*{Sources and Further Reading}
The following reputable organisations provide further information on melasma:
\begin{itemize}
    \item Australasian College of Dermatologists (ACD). \textit{A-Z of skin conditions: Melasma.} https://www.dermcoll.edu.au/
    \item DermNet. \textit{Comprehensive clinical resource on the causes and treatment of melasma.} https://dermnetnz.org/
    \item NHS. \textit{Overview of melasma and skin pigmentation changes.} https://www.nhs.uk/
    \item Mayo Clinic. \textit{Understanding melasma, symptoms, and cosmetic options.} https://www.mayoclinic.org/
    \item American Academy of Dermatology (AAD). \textit{Dermatologist guides to managing melasma.} https://www.aad.org/
\end{itemize}
